% Section 2.5 (Ben's diagram request): which tier serves which request, per
% system. Six requests R1..R6 in every panel. Drawn from the systems' own text:
%   (a) Pie s3.1 "We relocate m layers of the KV cache to the CPU, leaving n-m
%       layers on the GPU"; s3.3 swap in and swap out through GPU buffers.
%   (b) DirectKV s6 all K,V "in CPU memory"; abstract "GPU kernels to directly
%       access CPU-resident KV cache"; s1 streamed through SMEM.
%   (c) HeteroKV, s4: residents in HBM, riders in host DDR streamed into a
%       k-slot window, one forward pass; the link budget is X T0 per step.
% DAK, which also splits by request (concurrent work), is discussed in s8, not
% drawn. Colours from figstyle: HBM teal, DDR orange, H2D blue, D2H purple.
\begin{figure*}[t]
\centering
\begin{tikzpicture}[
  req/.style={draw=#1D, fill=#1F, line width=0.5pt, minimum width=5.6mm,
              minimum height=4.2mm, inner sep=0pt, font=\sffamily\tiny},
  half/.style={draw=#1D, fill=#1F, line width=0.5pt, minimum width=5.6mm,
               minimum height=2.1mm, inner sep=0pt},
  wts/.style={draw=black!45, fill=black!10, minimum height=4.2mm,
              inner sep=1pt, font=\sffamily\tiny},
  fwd/.style={cmp, minimum height=4.6mm, inner sep=1pt, font=\sffamily\scriptsize},
  tier/.style={font=\sffamily\scriptsize\bfseries, anchor=east},
  note/.style={font=\sffamily\scriptsize, text=inkS, align=center},
]
% ---------------------------------------------------------------- layout
% panel x offsets 0, 6.0, 12.0 cm; rows: compute 3.2, HBM 2.2, link 1.2, DDR 0.2
\foreach \px/\ttl in {0/{(a) Layer swap via HBM buffer},
                      6.0/{(b) Zero-copy (DirectKV)},
                      12.0/{(c) Split by request (\sysname{})}} {
  \begin{scope}[xshift=\px cm]
    \fill[bgA, rounded corners=3pt] (-0.15,-0.25) rectangle (5.55,3.75);
    \node[hd, anchor=south] at (2.7,3.78) {\ttl};
    \node[tier] at (0.9,2.2) {HBM};
    \node[tier] at (0.9,0.2) {DDR};
    \node[tier, font=\sffamily\scriptsize] at (0.9,1.2) {link};
  \end{scope}
}
% ------------------------------------------------------------------ (a) Pie
\begin{scope}[xshift=0cm]
  \node[fwd, minimum width=4.3cm] (fa) at (3.15,3.2) {one forward pass, R1--R6};
  \node[wts, minimum width=1.0cm] at (1.55,2.2) {weights};
  \foreach \i in {1,...,6} {
    \node[half=hbm] (ah\i) at ({1.95+0.5*\i},2.31) {};
    \node[half=ddr] (ad\i) at ({1.95+0.5*\i},0.09) {};
    \node[font=\sffamily\tiny] at ({1.95+0.5*\i},-0.1) {R\i};
  }
  \node[plain, minimum width=1.1cm, minimum height=2.1mm, inner sep=0pt,
        font=\sffamily\tiny] (abuf) at (3.7,1.95) {};
  \node[font=\sffamily\tiny, anchor=east] at (3.13,1.95) {buffer};
  \draw[h2d] (3.45,0.25) -- (3.45,1.84);
  \draw[d2h] (3.95,1.84) -- (3.95,0.25);
  \node[font=\sffamily\tiny, text=d2hC, anchor=west] at (4.0,1.05) {swap out};
  \draw[flow] (3.15,2.44) -- (3.15,2.96);
  \node[note] at (2.7,-0.05) {};
\end{scope}
% ------------------------------------------------------------- (b) DirectKV
\begin{scope}[xshift=6.0cm]
  \node[fwd, minimum width=4.3cm] (fb) at (3.15,3.2) {one forward pass, R1--R6};
  \node[wts, minimum width=1.0cm] at (1.55,2.2) {weights};
  \node[draw=black!40, densely dashed, minimum width=2.9cm, minimum height=4.2mm,
        font=\sffamily\tiny, text=inkS] at (3.65,2.2) {no KV in HBM};
  \foreach \i in {1,...,6} {
    \node[req=ddr] (bd\i) at ({1.95+0.5*\i},0.2) {R\i};
  }
  \draw[h2d, line width=1.3pt] (4.9,0.45) .. controls (5.35,1.4) and (5.35,2.6)
    .. (4.95,2.95);
  \node[font=\sffamily\tiny, text=h2dC, anchor=east, align=right]
    at (5.0,1.5) {straight into\\SMEM};
\end{scope}
% -------------------------------------------------------------- (c) HeteroKV
\begin{scope}[xshift=12.0cm]
  \node[fwd, minimum width=4.3cm] (fc) at (3.15,3.2) {one forward pass, R1--R5};
  \node[wts, minimum width=1.0cm] at (1.55,2.2) {weights};
  \foreach \i in {1,2,3} {
    \node[req=hbm] (ch\i) at ({1.95+0.5*\i},2.2) {R\i};
  }
  \node[draw=hbmD, fill=white, line width=0.5pt, minimum width=1.05cm,
        minimum height=4.2mm, inner sep=0pt, font=\sffamily\tiny]
    (cwin) at (4.55,2.2) {$k$ slots};
  \foreach \i in {4,5} {
    \node[req=ddr] (cd\i) at ({2.95+0.5*(\i-4)+0.75},0.2) {R\i};
  }
  \node[draw=black!40, densely dashed, fill=white, minimum width=5.6mm,
        minimum height=4.2mm, inner sep=0pt, font=\sffamily\tiny, text=inkS]
    (cd6) at (5.05,0.2) {R6};
  \draw[h2d] (4.0,0.45) -- (4.3,1.98);
  \draw[h2d] (4.45,0.45) -- (4.6,1.98);
  \draw[decorate, decoration={brace, amplitude=2.5pt, mirror}, black!60]
    (3.4,0.75) -- (3.4,1.65);
  \node[font=\sffamily\tiny, anchor=east, align=right] at (3.32,1.2)
    {at most $\Xlink T_0$\\bytes per step};
  \draw[flow] (2.95,2.44) -- (2.95,2.96);
  \draw[flow] (4.55,2.44) -- (4.55,2.96);
  \node[font=\sffamily\tiny, text=inkS, anchor=north] at (5.05,-0.02) {waits};
\end{scope}
\end{tikzpicture}
\caption{\textbf{One path per byte, or both.} (a) Layer swapping~\cite{pie} keeps a
share of every request's layers, sized online, in host memory and swaps them through an HBM buffer, in
and back out. (b) DirectKV keeps all KV in host memory and streams it into the
kernel's shared memory; HBM holds no KV. (c) \sysname{} splits the batch by
request: residents (R1--R3) read HBM, riders (R4, R5) stream from DDR into a
$k$-slot window within the link's per-step budget, and both run in one forward
pass. R6 waits.}
\label{fig:tiers}
\end{figure*}
